\section{Experiments}

\paragraph{Runs.}
We report three bounded API experiments, a 50-scene \texttt{gpt-5.5} speaker-generation audit, and a selected-message cross-model audit. The mixed pilot contains 50 scenes from the four initial scenario families; the perspective-stress replication contains 50 fresh perspective-shift scenes; and the partial-observability support run contains 50 private-landmark scenes. In all runs, the speaker model generates candidates, local selectors choose final messages, and held-out API listeners evaluate those messages. The cross-model audit reuses the same speaker candidates and selected messages, then evaluates them with \texttt{gpt-4.1-nano}, \texttt{gpt-5.4-nano}, and \texttt{gpt-5.5} held-out listener prompts. We use paired scene-level comparisons for method differences.

\paragraph{Model versions and reproducibility.}
The cache records exact response model versions: \texttt{\footnotesize gpt-5.4-nano-2026-03-17} for the main API runs, \texttt{\footnotesize gpt-4.1-nano-2025-04-14} for the alternate-model audits, and \texttt{\footnotesize gpt-5.5-2026-04-23} for the frontier-listener and speaker follow-up. The repository includes generated prompt and schema documentation, a token-accounting report over 7,171 cached responses, a benchmark integrity audit that checks generated messages and listener choices, and an artifact guide mapping paper claims to source scenes, cached result files, tables, and verifier outputs. As a benchmark-scale artifact sanity check, the release also includes a no-API 600-scene balanced local sweep: mirror self-play has 1.00 same-play but 0.63 cross-play success, while population-play, template, and oracle candidate each reach 1.00.

\paragraph{What is being tested?}
The experiments do not test whether parameter updates improve pragmatic reasoning. Instead, they test a smaller but common mechanism: using interaction-derived success to select among possible messages. This mechanism appears in reranking, rejection sampling, self-refinement, and data filtering. If selection against one partner can inflate apparent success in a controlled reference game, then cross-play is a useful diagnostic before moving to more expensive training or richer social environments.

\section{Results}

\paragraph{Finding 1: same-play success overstates held-out communicative success.}
In the mixed pilot, mirror self-play obtains 1.00 same-play success but 0.91 cross-play success; population-play obtains 1.00 cross-play success on the same scenes, with a paired population-minus-mirror difference of 0.09. Scenario-level results localize the effect: unique-attribute and distractor-contrast scenes are at ceiling, relational-reference scenes show a smaller mirror gap, and perspective-shift scenes are the hard case. This rules out a uniform degradation story; the self-play gap appears where listener perspective and ambiguity matter.

\paragraph{Finding 2: the gap replicates and transfers across model families.}
In 50 fresh perspective-shift scenes, mirror self-play again achieves 1.00 same-play success but only 0.81 cross-play success. Population-play reaches 1.00 cross-play success, matching the oracle candidate upper bound, and the paired population-minus-mirror difference is 0.19. Under an alternate held-out model family, \texttt{gpt-4.1-nano}, direct and shortest baselines become weaker, mirror self-play drops to 0.71 cross-play success, and population-play again reaches 1.00; the paired population-minus-mirror difference is 0.29.

Table~\ref{tab:crossmodel} summarizes the broader cross-model audit. The same selected messages are evaluated under three held-out listener families on perspective stress and partial observability. The mirror self-play gap is not repaired by \texttt{gpt-5.5}; the population-minus-mirror difference is 0.33 on perspective stress and 0.35 on partial observability. Population-play reaches 1.00 in every row. For the selected-message \texttt{gpt-5.5} rows, this also establishes an oracle ceiling of 1.00 without a new all-candidate listener sweep, because the population-selected message is one of the candidates and already succeeds for every held-out listener. A scene-level overlap audit shows that 20 of 22 alternate-model perspective mirror-failure scenes and 26 of 26 partial-observability mirror-failure scenes also fail under \texttt{gpt-5.5}, and all \texttt{gpt-5.5} mirror failures are symbolic-verifier positives.

\begin{table}[H]
\centering
\small
% Begin inlined tables/cross_model_listener_audit
\begin{tabular}{llrrrrr}
\toprule
Setting & Listener & Direct & Mirror & Population & Oracle & Pop.-Mirror \\
\midrule
Perspective stress & gpt-5.4-nano & 0.71 & 0.81 & 1.00 & 1.00 & 0.19 \\
Perspective stress & gpt-4.1-nano & 0.55 & 0.71 & 1.00 & 1.00 & 0.29 \\
Perspective stress & gpt-5.5 & 0.79 & 0.67 & 1.00 & 1.00 & 0.33 \\
Partial observability & gpt-5.4-nano & 0.74 & 0.67 & 1.00 & 1.00 & 0.33 \\
Partial observability & gpt-4.1-nano & 0.70 & 0.66 & 1.00 & 1.00 & 0.34 \\
Partial observability & gpt-5.5 & 0.74 & 0.65 & 1.00 & 1.00 & 0.35 \\
\bottomrule
\end{tabular}
% End inlined tables/cross_model_listener_audit
\caption{Cross-model held-out listener audit. GPT-5.5 rows reuse the existing speaker candidates and selected messages; the oracle is 1.00 because population-play already selects an all-listener-successful candidate in every scene.}
\label{tab:crossmodel}
\end{table}

We also regenerate all 50 perspective-stress speaker candidate sets with \texttt{gpt-5.5}. Stronger generation improves direct first-candidate success to 0.99 and yields robust candidates by $K=2$ in every scene, but local mirror selection still reaches only 0.85 cross-play success despite 1.00 same-play success. Population-play again reaches 1.00, giving a paired population-minus-mirror difference of 0.15.

\paragraph{Finding 3: coordinate fallbacks are useful, but selector quality matters without them.}
The strongest population-play result uses candidate sets that include an explicit row/column fallback. Table~\ref{tab:nocoord} asks whether the result is merely a coordinate fallback effect by removing messages that explicitly mention row or column numbers and recomputing selection without new API calls. Robust non-coordinate candidates remain available: the no-coordinate oracle reaches 0.79, and consensus+info reaches 0.76. However, naive population-play falls to 0.42 because the local training listeners often assign equal or low scores to more informative non-coordinate descriptions, causing the length penalty to favor short ambiguous messages. The same cache-only filter on the mixed pilot gives a no-coordinate oracle of 0.96 and consensus+info success of 0.92.

\begin{table}[H]
\centering
% Begin inlined tables/perspective_stress50_gpt41nano_no_coord
\begin{tabular}{lrrrr}
\toprule
Method & Cross-play & Same-play & Gap & Tokens \\
\midrule
Direct first & 0.55 & - & - & 11.7 \\
Best-of-K shortest & 0.37 & - & - & 4.6 \\
Mirror self-play & 0.58 & 1.00 & 0.42 & 8.3 \\
Population-play & 0.42 & 0.52 & 0.10 & 5.0 \\
Consensus+info & 0.76 & - & - & 12.5 \\
Informativeness prior & 0.76 & - & - & 12.5 \\
Oracle candidate & 0.79 & - & - & 9.4 \\
\bottomrule
\end{tabular}
% End inlined tables/perspective_stress50_gpt41nano_no_coord
\caption{No-exact-coordinate ablation on the alternate-model audit. We remove candidate messages that explicitly mention row or column numbers, then recompute selection while reusing cached held-out listener evaluations.}
\label{tab:nocoord}
\end{table}

A candidate-role audit shows why the ablation behaves this way. In the full perspective-stress and partial-observability runs, population-play selects the coordinate-fallback slot in 1.00 of scenes; in the alternate-model perspective audit, this is 50 coordinate selections versus 17 for mirror self-play. Once exact coordinates are removed, naive population-play selects a spatially informative message in only 4 of 50 perspective-stress scenes and drops to 0.42 cross-play success. Consensus+info falls back to the informativeness prior in all 50 stress scenes, selecting spatially informative messages throughout and reaching 0.76. In the mixed no-coordinate pilot, consensus+info uses local consensus in 32 scenes and the informativeness fallback in 18 scenes.

The dedicated API K=8 no-coordinate run changes this interpretation from a limitation to a sharper diagnostic. A \texttt{gpt-5.5} speaker prompt that forbids exact row/column language produces 400 candidates over the 50 perspective-stress scenes; 400 of 400 generated candidates survive the exact-coordinate filter. The K=8 no-coordinate oracle is 1.00 and the shortest-candidate selector also reaches 1.00, so robust non-coordinate expressions are available. Increasing from the K=4 filtered baseline to K=8 improves mirror from 0.85 to 0.95 and population-play rises from 0.83 to 0.99, but consensus+info drops from 0.99 to 0.90. The safe conclusion is therefore not that one heuristic selector dominates; it is that cross-play exposes when selection, rather than generation, is the remaining bottleneck.

\begin{figure}[H]
\centering
% Begin inlined tables/crossplay_gap_bars
\begin{tabular}{lll}
\toprule
Setting & Gap & Same-play minus cross-play \\
\midrule
Mixed 50, mirror & 0.09 & \rule{0.93cm}{6pt} \\
Mixed 50, population & 0.00 &  \\
Perspective stress, mirror & 0.19 & \rule{1.87cm}{6pt} \\
Perspective stress, population & 0.00 &  \\
\bottomrule
\end{tabular}
% End inlined tables/crossplay_gap_bars
\caption{Cross-play gap for mirror and population selection. The gap is same-play success minus held-out cross-play success.}
\label{fig:gaps}
\end{figure}

\paragraph{Failure decomposition and alternative explanations.}
The oracle result is important: in the perspective-stress replication, the API speaker usually generated a robust message, and the failure is selection. A cache-only selection-regret audit makes this decomposition explicit. In the full-candidate API runs, population-play has zero regret in the mixed, perspective-stress, alternate-model, and partial-observability settings, while mirror self-play leaves regrets of 0.09, 0.19, 0.29, and 0.33 respectively. In the no-coordinate ablations, consensus+info reduces regret to 0.04 in the mixed pilot, 0.03 in the perspective-stress audit, and 0.01 in partial observability, showing that most remaining failures are selector failures rather than missing candidates.

An even stricter candidate-pool audit asks whether any candidate succeeds with every held-out listener in a scene. Every full-candidate run has at least one such robust candidate in all 50 scenes, and population-play selects a robust candidate in 1.00 of scenes. Mirror robust-selection rates are 0.86, 0.72, 0.56, and 0.50. After exact coordinates are removed, robust non-coordinate candidates remain available in 0.92, 0.66, and 1.00 of scenes for the mixed, alternate-model perspective, and partial-observability runs, and consensus+info matches the per-scene oracle in 0.94, 0.92, and 0.98 of scenes.

A uniform-random candidate baseline rules out the explanation that any candidate would do. Random expected success is 0.87, 0.72, 0.66, and 0.75 in the four full-candidate runs, so population-play improves over random by 0.14, 0.29, 0.34, and 0.25. In the no-coordinate runs, consensus+info improves over random by 0.10, 0.22, and 0.32. A message-length audit rules out a pure verbosity explanation: population-play messages average 6.82, 10.22, 10.22, and 9.66 tokens, and direct messages are longer than mirror in the perspective and partial-observability runs while remaining less successful.

A prefix-$K$ candidate-budget audit shows why the four-candidate protocol is useful. In the alternate-model perspective audit, robust-scene coverage rises from 0.38 at $K=1$ to 0.66 at $K=3$ and 1.00 at $K=4$; in partial observability, it rises from 0.62 at $K=1$ to 1.00 at $K=3$. A held-out listener-disagreement audit shows that mirror self-play yields split held-out outcomes in 0.08, 0.18, 0.36, and 0.50 of scenes across the full-candidate runs, while population-play has 0.00 split outcomes in all four. In no-coordinate partial observability, consensus+info reduces split outcomes from mirror's 0.54 to 0.02. A listener-confidence audit further shows that self-reported uncertainty is not enough: in the alternate-model perspective runs, ambiguity flags are 0.00 while mirror self-play still produces high-confidence failures.

\paragraph{API-listener leave-one-out mechanism check and Partial-observability support check.}
To check that the result is not only an artifact of deterministic local selectors, we also run a cache-only leave-one-listener-out analysis over the API all-candidate evaluations. For each scene, one API listener prompt is held out; API-mirror selection uses one different API listener prompt, while API-population selection uses the other two. Population selection improves over API-mirror selection in all four cached runs, with larger gains for the alternate-model perspective audit (0.94 vs. 0.85) and partial observability (0.94 vs. 0.77). Because this analysis reuses the cached API listener pool for post-hoc splits, we treat it as a mechanism check rather than as the main held-out evaluation.

As an additional support check, we run a 50-scene partial-observability stress setting in which the speaker sees a private landmark hidden from the listener. The prompt instructs the speaker to use only listener-visible information. In this setting, mirror self-play obtains 0.67 cross-play success despite 1.00 same-play success, while population-play and the oracle both reach 1.00. After removing exact row/column candidates, robust non-coordinate messages remain available: the no-coordinate oracle is 1.00 and consensus+info reaches 0.99. The candidate audit finds zero explicit private-landmark references, so the failures reflect under-informative messages that fit both the target and distractor rather than object-ID leakage or hidden-landmark leakage. Rubric coding supports this interpretation: all 50 full-run mirror failures and all 59 no-coordinate mirror failures are underspecified-distractor choices.
